在热力学第一定律建立以前，人们曾反复尝试制造一种不需要任何能量供应、却能持续对外做功的机器，
也就是所谓的第一类永动机。所有这类尝试的失败，使人们逐渐认识到：必定存在一条普遍规律，使这种
机器从根本上不可能实现。今天我们知道，这条规律就是能量的转化和守恒定律。它是十九世纪科学发展中
最重要的成就之一；热力学第一定律则是能量守恒在包含热现象的系统中的定量陈述。

对于一个既不与外界交换物质、也不交换能量的孤立系统，总能量始终保持不变；若系统与环境之间存在
物质或能量交换，那么特定过程前后系统总能量的改变量，必然等于净输入系统的总能量。这里的总能量
包括内能、整体动能及整体势能；流入、流出的物质也能携带能量，不能只计算热量和边界功。
这一定律把一种新的宏观能量形式——内能——同热量、功以及其他可能的能量交换方式联系起来。

\perpetualMotionFaithful
\diagramcaption{图 20.1：三类第一类永动机设想。转轮、循环配重和径向摆锤的结构不同，却都企图在没有净能量输入时持续输出功。}

图20.1中的装置并不是因为“某个齿轮比例没有调好”才失败。只要装置完成一个循环并回到原状态，内能
的净变化就是零；若又没有从外界吸收能量，便不可能有正的净功输出。后文将把这一判断写成可计算的
方程。

\subsection*{内能}

热力学中的内能（internal energy）定义为包含在给定系统中的总能量，扣除系统质心整体运动带来的
动能，以及系统在外场中的整体势能。内能由系统内部微观粒子的运动状态和相互作用决定。在最简单的
经典图像中，它可以理解为微观粒子热运动的动能与粒子间相互作用势能之和。

一般情形下，系统内能的绝对值很难直接确定，因为分子之间存在复杂的相互作用。幸运的是，在经典物理
中，真正进入可观测规律的通常是内能的改变 $\Delta U$。这与海面航行的类比相似：许多问题并不需要
知道海水总量或海底的绝对深度，只需要知道海面高度在过程中的变化。内能的转化过程仍由底层的动力学
规律支配，但热力学把大量微观自由度压缩为少数宏观状态量。

只有在某些简单模型中，内能可以直接写出。例如单原子理想气体：构成气体的原子在所研究的过程中内部
结构不变，气体处于绝对温度 $T$ 的热平衡状态，分子数为 $N$，物质的量为 $n=N/N_A$。理想气体忽略
分子间相互作用，因而不存在分子间势能；每个单原子分子有三个平动自由度。由能量均分定理，每个分子
平均具有 $3k_BT/2$ 的平动动能，所以全部分子的内能为
\begin{equation}
  \boxed{U=\frac32Nk_BT=\frac32nRT}. \tag{20.1}
\end{equation}
由此可见，在热平衡下，一定量单原子理想气体的内能只与温度有关。

以氧气、氮气为代表的双原子气体稍复杂。除质心的三个平动自由度外，分子还可以转动，也可能沿两原子
连线方向振动。若不加判断地使用经典能量均分定理，会得到高温极限下的 $U=7nRT/2$；实际测量却显示，
常温下振动自由度并没有充分激发，气体的行为更像只有平动和转动自由度。这个问题长期困扰经典物理，
直到量子理论建立后，人们才理解：相应能级间隔若远大于 $k_BT$，该自由度会被冻结；温度升高后，
振动贡献才逐渐显现。

对于不能忽略分子间势能的物质，内能通常不仅与温度有关，也会与体积有关，因为体积改变会改变分子
间距和相互作用势能。尽管内能的具体表达式可能很复杂，但有一点极其重要：当系统的宏观平衡状态完全
给定时，其内能也随之确定。内能是状态函数；系统以何种过程来到同一个状态，不会改变该状态的内能。

\subsection*{热交换：传导、对流与辐射}

当系统与外界环境存在温度差时，就具备发生热量交换的条件。约定系统从状态1到状态2的过程中，从外界
净输入热量记作 $Q_{1\to2}$。若 $Q>0$，能量以热的形式净流入系统；若
$Q<0$，系统净放热。后文第一定律中的 $Q^{\rm in}$ 也表示这一带符号的净热量，即 $Q^{\rm in}=Q$。
计算循环效率时，另把各吸热段的热量总和记为正量 $Q_h$，各放热段排出的热量总和记为正量 $Q_c$，则
\begin{equation*}
  Q_{\rm cyc}=Q_h-Q_c.
\end{equation*}
这两个正的累计热量一般不相等，也不是互为负数。
热量像功一样，是过程量而不是状态量：不能说“某个状态含有多少热量”，只能说某一过程中传递了多少
热量。

温度差只创造了热交换的条件，能量真正传递还需要具体机制。常见机制包括直接接触传导、流体对流和
热辐射。两个温度不同的固体直接接触，或通过某种介质接触时，热端粒子的平均能量较高；相互作用使
能量在大量碰撞中向冷端传递。宏观上热端降温、冷端升温，这就是热传导。在最简单的近似下，热流率
与温差成正比，比例系数还取决于材料和几何尺寸，所以不同介质的传热能力差别显著。

\redrawimage[0.62\linewidth]{figure20-2-redraw-v3.png}
\diagramcaption{图 20.2：手掌附近的对流羽流示意。空气受热、密度减小并上升；颜色表示相对冷热，不代表实测温度。}
图20.2中，手掌是局部热源；上方羽流逐渐扩展，表明热空气上升时还会卷入周围空气。

气体只靠分子扩散传导热量的能力通常不强；当介质是重力场中的流体，并在竖直方向存在足够大的温度梯度
时，会出现效率更高的对流。以地表附近的空气为例，随高度增加，大气压强和密度逐渐减小。若低处一团
较热气体受到扰动向上运动，到达高处后密度仍小于周围气体，它将继续受到向上的浮力，并把所携带的
热量送往高处。对流一旦形成，流体的整体运动会迅速交换各区域的能量，趋向减小温差。

\stellarInteriorRevision
\diagramcaption{图 20.3：太阳与红巨星的内部剖面及表面对流细节。箭头表示对流区物质的循环运动；两星大小与内部层厚不按同一比例。}
图20.3左侧太阳的核心产生能量，能量先穿过辐射区，再经外层对流区向外传递；右侧红巨星具有扩展的
对流包层和相对紧凑的内部区域。下方局部图中，亮区对应较热物质上涌，暗边界对应较冷物质下沉。
这说明对流不仅发生在地球空气中，也能成为恒星内部的重要输能方式；不能把表面纹理直接当作内部结构。

传导和对流都需要物质介质，热辐射则可以穿过真空。太阳与地球之间的能量交换就是典型例子。所有温度
高于绝对零度的物体都会发出电磁辐射，辐射带走能量；温度越高，单位面积的总辐射功率越大。不同物质
的发射能力不同。吸收能力越强的表面，在同一温度下通常也具有更强的发射能力。能够吸收全部入射辐射
的理想物体称为黑体（black body）。黑体温度为 $T$ 时，单位表面积的总辐射功率为
\begin{equation}
  \boxed{j_{\rm rad}=\sigma T^4}, \tag{20.2}
\end{equation}
其中 $\sigma\approx\SI{5.670374419e-8}{W.m^{-2}.K^{-4}}$ 是 Stefan--Boltzmann 常数。四次方关系
说明，温度升高时热辐射能力会急剧增强。

\redrawimage[0.72\linewidth]{figure20-4-redraw-v2.png}
\diagramcaption{图 20.4：手部热辐射强弱示意。掌心、指尖、手腕和边缘的温度不同；图中颜色不是实验测量值。}
图20.4的各部位即使没有接触远处探测器，也能通过电磁辐射交换能量。

\heatTransferMapDiagram
\figcap{热传递与边界功的路径。热辐射不需要物质介质，功对应边界受力和位移。}\label{fig:heat-transfer-map}
图\ref{fig:heat-transfer-map}中的 $Q$ 与 $W$ 都跨越系统边界，却对应不同的能量交换机制。

上述三种热交换都由温度差驱动，但路径不同：传导靠微观碰撞逐步传能，对流还伴随流体团块的宏观位移，
辐射则以电磁波携带能量。判断题目中的热传递方式时，不能只看“有没有接触”，还要看是否发生整体流动、
是否允许真空传播，以及过程所经历的时间尺度。

\subsection*{机械功与准静态过程}

当系统与环境之间存在压强差，并且系统边界允许移动时，就会发生另一类能量交换——机械功。系统膨胀
时可对环境做功；系统被压缩时，环境对系统做功。

设圆柱形气缸横截面积为 $A$，封闭气体压强为 $P$。气体作用在活塞上的力为 $F=PA$。活塞向外移动
极小距离 $\dd y$ 时，体积增量 $\dd V=A\dd y$，气体对外所做的元功为
\begin{equation*}
  \dd W^{\rm out}=F\dd y=P\,\dd V.
\end{equation*}
这一结果与容器的具体形状无关。若可逆过程中压强随体积按 $P(V)$ 改变，从 $V_i$ 变到 $V_f$ 的对外功
就是 $P$--$V$ 图像下的有向面积：
\begin{equation}
  \boxed{W^{\rm out}_{i\to f}=\int_{V_i}^{V_f}P(V)\,\dd V}. \tag{20.3}
\end{equation}

\pistonWorkFaithful
\diagramcaption{图 20.5：活塞发生微小位移时，气体体积增加 $\dd V=A\dd y$；同一过程在 $P$--$V$ 图中曲线下的面积就是气体对外功。}
图20.5右图的 $i,f$ 与横轴上的 $V_i,V_f$ 一一对应，阴影上边界就是所积分的同一条压强曲线。

若改用环境对系统所做的功，符号相反：$\dd W^{\rm in}=-P\dd V$。但把边界处的功直接写成
$P\dd V$ 隐含一个条件：外界压强与系统压强在全过程中几乎相等，只保留足以推动边界移动的无穷小差别。
因此过程必须足够缓慢，使系统在每一时刻都接近平衡态；这样的理想过程称为准静态过程。若再忽略摩擦、
黏滞、有限温差传热等耗散，准静态过程可进一步视为可逆过程。

真实压缩或膨胀总会偏离这一理想极限，但竞赛题若给出“缓慢”“光滑”“无摩擦”等条件，通常正是在
授权使用准静态、可逆近似。非准静态过程仍满足能量守恒，只是系统内部可能没有统一压强，不能直接用
系统平衡压强代入式(20.3)计算边界功。

\subsection*{热力学第一定律}

热交换和机械功是系统与环境交换能量的两种最常见方式。本章主要讨论不交换物质的封闭系统。
一般应写 $\Delta(U+E_k+E_p)=Q+W^{\rm in}$。若系统从状态1经任意过程到达状态2，
整体动能和外场势能的变化可忽略，只考虑热量和功两种交换形式，取“流入系统为正”，则可简化为
\begin{equation}
  \boxed{U_2-U_1=Q^{\rm in}_{1\to2}+W^{\rm in}_{1\to2}}. \tag{20.4}
\end{equation}
若改用气体对外做功 $W^{\rm out}=-W^{\rm in}$，同一规律等价地写成
\begin{equation*}
  \boxed{Q^{\rm in}=\Delta U+W^{\rm out}}.
\end{equation*}
后一个形式在气体过程题中更常用：吸收的热量一部分增加内能，另一部分转化为气体对外功。

热力学第一定律是能量守恒的直接体现，因此对可逆、不可逆过程都成立。可逆性的作用不是决定能量是否
守恒，而是决定能否仅凭一条平衡态路径计算功和热量。内能 $U$ 是状态函数，$\Delta U$ 只取决于始末
状态；热量 $Q$ 与功 $W$ 是过程量，同样的两个状态之间，沿不同路径的数值可以不同。

若系统经历一个循环又回到初态，则 $\Delta U_{\rm cyc}=0$。因此
\begin{equation*}
  0=W^{\rm in}_{\rm cyc}+Q^{\rm in}_{\rm cyc},
  \qquad
  \boxed{W^{\rm out}_{\rm cyc}=Q^{\rm in}_{\rm cyc}}.
\end{equation*}
整个循环对外的净功必等于同一循环从外界吸收的净热量。若循环不从外部吸收任何能量，净功只能为零；
任何企图无能量输入而持续对外做功的第一类永动机都不可能存在。

\begin{warningbox}{先统一功的正方向}
有的教材把“气体对外做功”记为正，有的教材把“外界对气体做功”记为正。两种约定都可以，但不能在
同一推导中混用。本讲解答统一采用 $Q>0$ 表示气体吸热，$W^{\rm out}>0$ 表示气体对外做功，故
$Q=\Delta U+W^{\rm out}$。
\end{warningbox}

\subsection*{理想气体的等容、等压与等温过程}

热力学第一定律配合理想气体状态方程，可以进一步研究具体过程。理想气体的内能只与温度有关；在温区
不太宽、摩尔热容可近似为常量时，$n$ 摩尔气体有
\begin{equation*}
  \dd U=n c_v\,\dd T,
\end{equation*}
其中 $c_v$ 是定容摩尔热容。微分式也适用于随温度变化的 $c_v(T)$；但有限温差时一般应算
\[
 \Delta U=n\int_{T_1}^{T_2}c_v(T)\,\dd T.
\]
只有在该温区内 $c_v$ 近似不变，才可把它提出积分得到 $n c_v\Delta T$。

理想气体经历等容过程时，$\dd V=0$，因而不发生体积功。吸收的热量全部表现为内能增加：
\begin{equation*}
  \boxed{Q_V=n c_v\Delta T}.
\end{equation*}
单原子理想气体在经典适用范围内 $c_v=3R/2$。双原子气体的热容随温度改变：低温时部分转动自由度
冻结，常温时主要有平动和转动贡献，高温时振动自由度逐渐激发。

\heatCapacityRedraw
\diagramcaption{图 20.6：双原子气体热容的自由度激发示意。三个平台为 $3R/2,5R/2,7R/2$，横轴各十倍温区等宽；不计解离、电子激发和相变。}
图20.6中每增加一个平台，表示又有一组自由度开始贡献热容，不是气体发生了不连续的温度跳变。

等压过程比等容过程多了膨胀功。温度增加 $\Delta T$ 时，理想气体状态方程给出
$P\Delta V=nR\Delta T$，于是
\begin{equation*}
  Q_P=n c_v\Delta T+P\Delta V
  =n(c_v+R)\Delta T=n c_p\Delta T.
\end{equation*}
因此理想气体满足 Mayer 关系
\begin{equation*}
  \boxed{c_p=c_v+R}.
\end{equation*}
$c_p$ 是一摩尔气体在等压过程中升高一开尔文所需的热量。它比 $c_v$ 大一个 $R$，正是因为等压升温
时还要为膨胀功提供能量。

若气体全过程同恒温热库接触并始终保持温度 $T$，便经历等温过程。对理想气体，$U$ 只与 $T$ 有关，
所以等温过程 $\Delta U=0$。气体等温膨胀时吸收的热量完全转化为对外功。$n$ 摩尔理想气体从 $V_1$
等温膨胀到 $V_2$ 时
\begin{equation}
  \boxed{Q^{\rm in}=W^{\rm out}
  =\int_{V_1}^{V_2}P\,\dd V
  =nRT\ln\!\frac{V_2}{V_1}}. \tag{20.5}
\end{equation}
理想气体具有“等温即内能不变”的简单性质。真实物质的内能通常还依赖体积，不能把这条结论直接
推广到液体、固体或高密度真实气体。

\processComparisonDiagram
\figcap{四类常用气体过程的约束。等压热量式和绝热幂律分别使用常热容与可逆绝热近似。}\label{fig:process-constraints}
图\ref{fig:process-constraints}应按已知条件逐项选用；例如 $Q=0$ 本身并不足以推出 $PV^\gamma$ 恒定。

\subsubsection*{同一对初末态，不一定有同样的功和热}

以内能只依赖温度的理想气体为例，设初态为 $(T_0,V_1)$，末态为 $(T_0,V_2)$，且 $V_2>V_1$。
若气体接触温度为 $T_0$ 的热库，作可逆等温膨胀，则
\[
 \Delta U=0,\qquad W^{\rm out}=nRT_0\ln\frac{V_2}{V_1},\qquad Q=W^{\rm out}>0.
\]
若改在绝热刚性容器中撤去隔板，让气体向真空自由膨胀到同一体积，则外压为零，
整个容器外壁也不移动，所以 $W^{\rm out}=0$；绝热给出 $Q=0$，于是仍有 $\Delta U=0$。
理想气体的平衡末温仍为 $T_0$，两条路径连接相同的初末态，但功、热不同。
自由膨胀的中间过程并非一串均匀平衡态，不能把平衡物态方程 $p=nRT_0/V$
当作边界上的实际外压去积分。写体积功之前必须说明使用哪一侧的压强，以及边界是否准静态。

同理，$Q=0$ 不能单独推出 $TV^{\gamma-1}$ 恒定；刚才的自由膨胀就是反例。
内能的改变量由状态决定，热和功则由实际能量交换的方式决定，这是后续所有过程题选式的依据。

\subsection*{绝热过程}

温差出现后，真正传递有限热量需要一定时间。现实中的快速过程虽然可能存在明显温差，但在过程持续的
短时间内，热交换相对于功和内能改变很小，可以近似为绝热过程。绝热过程定义为系统与环境之间没有
热量交换，即 $Q=0$。若没有其他能量交换，绝热膨胀对外做的功完全来自内能降低；绝热压缩时，外界
做功则转化为内能增加。
快速并不是绝热的定义：隔热良好的系统也可以缓慢地发生绝热变化。若要同时使用“绝热”和“准静态”
近似，过程时间应长于系统内部恢复力学平衡的时间，又短于显著漏热的时间。

考虑 $n$ 摩尔理想气体从平衡状态 $(P,V,T)$ 作无穷小可逆绝热变化。由
\begin{equation*}
  T=\frac{PV}{nR},
  \qquad
  \dd T=\frac{P\dd V+V\dd P}{nR},
\end{equation*}
以及第一定律 $\dd U+P\dd V=n c_v\dd T+P\dd V=0$，先代入温度微分：
\[
 \frac{c_v}{R}(P\dd V+V\dd P)+P\dd V=0,
 \qquad (c_v+R)P\dd V+c_v V\dd P=0.
\]
利用 $c_p=c_v+R$，再除以 $c_vPV$，得到
\begin{equation}
  \boxed{\gamma\frac{\dd V}{V}+\frac{\dd P}{P}=0},
  \qquad
  \boxed{\gamma=\frac{c_p}{c_v}}. \tag{20.6}
\end{equation}
$\gamma$ 称为绝热指数。以下幂律公式还要求 $c_v,c_p$ 在所研究温区内为常数。对微分式从初态积分至末态：
\[
 \int_{P_1}^{P_2}\frac{\dd P}{P}+\gamma\int_{V_1}^{V_2}\frac{\dd V}{V}=0,
 \qquad \ln\frac{P_2}{P_1}+\gamma\ln\frac{V_2}{V_1}=0.
\]
把对数相加并取指数，得 $P_2V_2^\gamma=P_1V_1^\gamma$，即泊松（Poisson）绝热方程
\begin{equation}
  \boxed{PV^\gamma=\mathrm{const}}. \tag{20.7}
\end{equation}

利用 $P=nRT/V$，有 $PV^\gamma=nRTV^{\gamma-1}$；定量气体的 $nR$ 不变，故
\begin{equation}
  \boxed{TV^{\gamma-1}=\mathrm{const}}. \tag{20.8}
\end{equation}
这一形式适合研究封闭气体在热机中的体积、温度变化。若代入 $V=nRT/P$，则
$PV^\gamma=(nR)^\gamma T^\gamma/P^{\gamma-1}$，所以
\begin{equation}
  \boxed{\frac{T^\gamma}{P^{\gamma-1}}=\mathrm{const}}. \tag{20.9}
\end{equation}
压强和温度都是强度量，式(20.9)不显含气体总量，适合跟踪开放环境中不与外界混合的
同一小气团，并不意味着可随意用于有质量流入、流出的系统。再以密度
$\rho$ 代替体积，可写成
\begin{equation}
  \boxed{P=\mathrm{const}\cdot\rho^\gamma}. \tag{20.10}
\end{equation}
压强与密度的关系是描述气体物态及其变化的直接工具。
这里每条绝热线的常数取决于气体状态。对固定物质的量、常热容理想气体，下一章的熵公式给出
\[
 S_2-S_1=n c_v\ln\frac{P_2V_2^\gamma}{P_1V_1^\gamma}.
\]
所以同一条可逆绝热线也是等熵线；这个常数不是焓。若热容明显随温度变化，须改用
$\int_{T_1}^{T_2}c_v(T)\dd T/T+R\ln(V_2/V_1)=0$，不能套固定指数的幂律。

\adiabaticCurvesDiagram
\figcap{从同一初态出发，理想气体等温线与可逆绝热线的比较。}\label{fig:adiabatic-comparison}
图\ref{fig:adiabatic-comparison}中，在任一相同的末体积处，下支的压强较低，对应膨胀降温的绝热过程。

在相同初态出发的膨胀过程中，绝热线比等温线下降得更陡。等温膨胀可以从热库持续吸热来补偿对外功；
绝热膨胀没有热量补充，只能消耗内能，温度随体积增大而下降。这也是水汽和空气团快速上升时降温的
基本原因。

\subsection*{重力场中形成对流的条件}

绝热方程的一个重要应用，是判断重力场中的理想气体能否形成对流。受扰动气团来到新高度后，会随外界
压强变化而膨胀或压缩。大气中高处压强较小，气团上升时迅速膨胀；若上升足够快，气团与周围空气来不及
交换显著热量，便可近似为绝热过程。

设大气具有已知的稳定温度分布 $T(h)$。在高度 $h$ 与 $h+\dd h$ 之间，静力平衡给出
\begin{equation*}
  \dd P=-\rho g\,\dd h=-\frac{\mu gP}{RT}\,\dd h,
\end{equation*}
其中 $\mu$ 为气体平均摩尔质量。把同一小气团从 $h$ 绝热抬升到 $h+\dd h$，由式(20.9)可得它的
温度改变
\begin{equation*}
  \dd T_1=\frac{\gamma-1}{\gamma}\frac{T}{P}\dd P
  =-\frac{\gamma-1}{\gamma}\frac{\mu g}{R}\dd h
  =-\frac{\mu g}{c_p}\dd h.
\end{equation*}
而环境中原本位于 $h+\dd h$ 处的空气温度为
\begin{equation*}
  T_2=T+\frac{\dd T}{\dd h}\dd h.
\end{equation*}
两团气体在同一高度具有近似相同压强和成分；温度较高者密度较小。要使受扰动上升后的气团比周围空气
更热、更轻，从而继续上升，必须满足
\begin{equation}
  \boxed{\frac{\dd T}{\dd h}<-\frac{\mu g}{c_p}}. \tag{20.11}
\end{equation}
因此，只要求低处比高处热还不够；环境温度随高度降低的速率必须大于气团的绝热降温率。

\pressureLayerRedraw
\diagramcaption{图 20.7：取横截面积为 $A$、厚度为 $\dd h$ 的水平气层，由竖直方向受力平衡得到 $\dd P=-\rho g\dd h$。}
图20.7中上、下表面压强差支撑薄层的重量；$\dd h>0$ 时 $\dd P<0$。

\convectionCriterionDiagram
\figcap{受扰动气团与环境温度的比较。两者来到同一高度后，压强近似相等。}\label{fig:convection-displacement}
图\ref{fig:convection-displacement}中，应比较上方气团的 $T+\dd T_1$ 与该高度环境的 $T+\dd T_2$，
不能只比较气团上升前后的温度。

\atmosphereProfileFaithful
\diagramcaption{图 20.8：地球大气的代表性温度、压强与分层。下方列出相关自然情景和飞行器，二者不是同一比例的图。}
图20.8上方的温度和压强对应同一个高度刻度；下方情景只帮助辨认物体，并不按高度比例排布。

从地表向上，对流层的平均温度通常随高度降低，但“随高度降低”本身不保证干对流不稳定。
例如约 $6.5\,\mathrm{K/km}$ 的平均温度递减率，小于干空气约 $9.8\,\mathrm{K/km}$ 的绝热递减率；
地表局部强加热或水汽凝结才可能使情况改变。再向上进入平流层，
温度可能随高度增加，对流受到抑制。云层多位于对流活跃区域，民航巡航高度则常接近平流层下部，以避开
低层大气中较强的天气扰动。水汽、二氧化碳等吸收地表长波辐射；平流层臭氧吸收太阳紫外线，
热层稀薄气体吸收更高能辐射。极光主要来自带电粒子激发高层气体，并不是空气热辐射的直接显示。
真实大气还受到水汽凝结、辐射吸收和地表加热等影响，式(20.11)给出的是
干燥理想气体的基本判据。

\setcounter{equation}{11}
\Needspace{10\baselineskip}
\begin{origexamplebox}{例题 20.1【正方形循环】}
小明制作了一个气缸，封住了一些理想气体，并让气体在 $P$--$V$ 图中完成三种顺时针循环。三种循环
具有相同的边界压强 $P_1,P_2$ 和边界体积 $V_1,V_2$：

（1）矩形循环；（2）上三角循环；（3）下三角循环；（4）比较这三种循环以及不同大小循环的效率。

下图给出三种循环的形状与运行方向。设气体的定容摩尔
热容为 $c_v$，并在所考察温区内近似不变。
\threeCyclesDiagram
\figcap{矩形、上三角与下三角循环。三个图使用相同的 $P_1,P_2,V_1,V_2$。}\label{fig:three-geometric-cycles}
\end{origexamplebox}

\begin{xesolution}
记
\begin{equation*}
  \Delta P=P_2-P_1,\qquad \Delta V=V_2-V_1,
  \qquad \alpha=\frac{c_v}{R}.
\end{equation*}
对任意过程，理想气体的温度满足 $nRT=PV$，所以
\begin{equation*}
  \Delta U=n c_v\Delta T=\alpha\Delta(PV),
  \qquad Q=\Delta U+W^{\rm out}.
\end{equation*}

图\ref{fig:three-geometric-cycles}中的箭头均围出正的循环功，但各段吸热量并不相同。

（1）矩形循环的净功就是矩形面积：
\begin{equation*}
  W_{\text{矩}}=\Delta P\,\Delta V.
\end{equation*}
吸热发生在左侧等容升压段与上侧等压膨胀段。两段吸热量分别为
\begin{equation*}
  Q_{V_1}=\alpha V_1\Delta P,
  \qquad
  Q_{P_2}=(\alpha+1)P_2\Delta V.
\end{equation*}
因此
\begin{equation}
  \boxed{\eta_{\text{矩}}=
  \frac{\Delta P\Delta V}
  {\alpha V_1\Delta P+(\alpha+1)P_2\Delta V}}.
\end{equation}

（2）上三角循环的顺时针面积为矩形的一半。它的吸热过程仍是完整的左侧等容升压段和上侧等压膨胀段；
从右上角沿斜线回到左下角时，$P,V,T$ 都下降，同时外界对气体做功，该段为放热。因此
\begin{equation}
  \boxed{\eta_{\rm upper}=
  \frac{\tfrac12\Delta P\Delta V}
  {\alpha V_1\Delta P+(\alpha+1)P_2\Delta V}}.
\end{equation}

（3）下三角循环从左下角沿直线膨胀到右上角时吸热，其余等容降压、等压压缩两段都放热。直线段对外功
等于梯形面积，故吸热量为
\begin{align*}
  Q_{\rm diag}
  &=\alpha(P_2V_2-P_1V_1)
    +\frac{P_1+P_2}{2}(V_2-V_1).
\end{align*}
于是
\begin{equation}
  \boxed{\eta_{\rm lower}=
  \frac{\tfrac12\Delta P\Delta V}
  {\alpha(P_2V_2-P_1V_1)+\tfrac12(P_1+P_2)\Delta V}}.
\end{equation}

为了比较，令矩形循环的吸热量为 $Q_{\text{矩}}$、净功为 $W_{\text{矩}}$。下三角斜边吸热量可化为
\begin{equation*}
  Q_{\rm diag}=Q_{\text{矩}}-\frac12W_{\text{矩}}.
\end{equation*}
又因为这里的矩形热机有 $0<W_{\text{矩}}<Q_{\text{矩}}$，所以
\begin{equation*}
  \boxed{\eta_{\text{矩}}>\eta_{\rm lower}>\eta_{\rm upper}}.
\end{equation*}

（4）循环画得“面积更大”并不自动意味着效率更高。效率只依赖无量纲比值
$P_2/P_1$、$V_2/V_1$ 和 $c_v/R$。若把循环关于原点相似放大，使这两个比值不变，功和吸热量按同一
比例增大，效率不变；若固定 $P_1,V_1$ 而改变 $P_2,V_2$，则必须代入上述公式具体比较。功率、每循环
输出功和效率是三个不同概念。
\end{xesolution}

\Needspace{10\baselineskip}
\begin{origexamplebox}{例题 20.2【新式气缸】}
小明制作了一种新式气缸，气体在其中完成下列循环：

（1）两绝热、两等温循环；（2）两绝热、两等体循环；（3）两绝热、两等压循环。

试用各状态点的温度、体积或压强，以最简单形式表示循环效率。取定量、定组成的理想气体，
摩尔热容在整个温区内近似不变，所有绝热段均可逆。下图依次为卡诺、奥托和布雷顿循环。
\standardHeatCyclesDiagram
\figcap{三种循环的状态点与方向。均按 $1\to2\to3\to4\to1$ 运行；各曲线由同一气体的状态方程确定。}\label{fig:three-standard-cycles}
\end{origexamplebox}

\begin{xesolution}
三种循环都用绝热段连接吸热、放热过程。绝热段本身 $Q=0$，关键是分别计算另外两段的吸、放热量。

在图\ref{fig:three-standard-cycles}中，卡诺循环的 $1\to2$ 为高温等温膨胀；奥托与布雷顿循环的
$1\to2$ 则为绝热压缩。下面的温度下标严格沿用各图的状态号。

（1）两条等温线温度分别为 $T_h,T_c$。高温等温膨胀吸热，低温等温压缩放热。由两条绝热线
$TV^{\gamma-1}=\mathrm{const}$ 有
\[
 T_hV_2^{\gamma-1}=T_cV_3^{\gamma-1},\qquad
 T_hV_1^{\gamma-1}=T_cV_4^{\gamma-1}.
\]
相除得 $(V_2/V_1)^{\gamma-1}=(V_3/V_4)^{\gamma-1}$。
记两边的体积比为 $r>1$，两等温段的吸、放热量大小因此为
\begin{equation*}
  Q_h=nRT_h\ln r,\qquad Q_c=nRT_c\ln r.
\end{equation*}
故 Carnot 循环效率为
\begin{equation}
  \boxed{\eta_{\rm C}=1-\frac{T_c}{T_h}}.
\end{equation}

（2）两条等容线位于 $V_{\min}$、$V_{\max}$，定义压缩比
$r=V_{\max}/V_{\min}>1$。等容吸、放热量为
\begin{equation*}
  Q_h=n c_v(T_3-T_2),\qquad Q_c=n c_v(T_4-T_1).
\end{equation*}
两条绝热线给出
\begin{equation*}
  \frac{T_2}{T_1}=\frac{T_3}{T_4}=r^{\gamma-1},
\end{equation*}
分别写成 $T_2=r^{\gamma-1}T_1$ 和 $T_3=r^{\gamma-1}T_4$，相减得
$T_3-T_2=r^{\gamma-1}(T_4-T_1)$。代入热量比并约去 $nc_v$，
得到 $Q_c/Q_h=1/r^{\gamma-1}$，奥托循环效率为
\begin{equation}
  \boxed{\eta_{\rm O}=1-\frac{1}{r^{\gamma-1}}}.
\end{equation}

（3）两条等压线压强分别为 $P_h,P_l$，定义压强比 $r_p=P_h/P_l>1$。等压吸、放热量为
\begin{equation*}
  Q_h=n c_p(T_3-T_2),\qquad Q_c=n c_p(T_4-T_1).
\end{equation*}
由绝热关系 $T^\gamma/P^{\gamma-1}=\mathrm{const}$，
\begin{equation*}
  \frac{T_2}{T_1}=\frac{T_3}{T_4}=r_p^{(\gamma-1)/\gamma}.
\end{equation*}
令 $b=r_p^{(\gamma-1)/\gamma}$，则 $T_2=bT_1,T_3=bT_4$，所以
$T_3-T_2=b(T_4-T_1)$。约去热量比中的 $nc_p$，得到布雷顿循环效率
\begin{equation}
  \boxed{\eta_{\rm B}=1-\frac{1}{r_p^{(\gamma-1)/\gamma}}}.
\end{equation}
三式的共同结构都是 $\eta=1-Q_c/Q_h$；区别在于吸、放热段受何种状态量约束。
这里 $Q_h,Q_c$ 均为正的热量大小。效率无量纲，且在温差、压缩比或压强比
分别趋于零、1、1时趋于零；同一气体、同一可逆绝热条件必须贯穿整个推导。
\end{xesolution}

\Needspace{10\baselineskip}
\begin{origexamplebox}{例题 20.3【变化的热容】}
小明发现，气体在一段过程中对应的过程热容未必保持为常数。为验证这一点，他设计了一摩尔理想气体
经历的过程
\begin{equation*}
  \frac{P}{P_0}+\frac{V}{V_0}=1.
\end{equation*}
设气体的定容摩尔热容为 $c_v$，求整个过程中摩尔过程热容随体积的变化关系。
\end{origexamplebox}

\begin{xesolution}
过程热容定义为 $C=\delta Q/\dd T$。它是过程属性，不等同于气体本身的 $c_v$ 或 $c_p$。
由 $P>0,V>0$，本题的物理区间为 $0<V<V_0$；两截距只是边界，不是气体可取的常温有限状态。由过程方程
\begin{equation*}
  P=P_0\left(1-\frac{V}{V_0}\right),
\end{equation*}
一摩尔理想气体的温度为
\begin{equation*}
  T(V)=\frac{PV}{R}
  =\frac{P_0}{R}\left(V-\frac{V^2}{V_0}\right),
\end{equation*}
从而
\begin{equation*}
  \frac{\dd T}{\dd V}=\frac{P_0}{R}\left(1-\frac{2V}{V_0}\right).
\end{equation*}

\variableHeatCapacityDiagram
\figcap{线性 $P(V)$ 与相应的 $T(V)$。温度抛物线的顶点严格位于两个体积截距的中间。}\label{fig:variable-heat}
图\ref{fig:variable-heat}中左侧压强随体积线性下降，右侧温度却先升后降，这正是过程热容变号的原因。

由第一定律
\begin{equation*}
  \delta Q=c_v\dd T+P\dd V,
\end{equation*}
所以
\begin{align*}
  C(V)
  &=c_v+P\frac{\dd V}{\dd T}\\
  &=c_v+R\frac{1-V/V_0}{1-2V/V_0}.
\end{align*}
即
\begin{equation}
  \boxed{C(V)=c_v+R\frac{V_0-V}{V_0-2V}}.
\end{equation}
在 $V=V_0/2$ 处，温度达到最大值，温度的一阶变化为零，但热量的一阶变化仍有
$\delta Q=P\dd V=(P_0/2)\dd V\ne0$，所以以 $\delta Q/\dd T$
定义的过程热容发散。发散不是气体储存了无穷能量，而是“用温度变化量作分母”在温度转折点失去良好
定义。令 $z=V/V_0$、$\alpha=c_v/R>0$，则
\[
 \frac{C}{R}=\frac{(\alpha+1)-(2\alpha+1)z}{1-2z},\qquad
 z_* =\frac{\alpha+1}{2\alpha+1}\in(1/2,1).
\]
因此 $0<z<1/2$ 时 $C>0$；$1/2<z<z_*$ 时 $C<0$；$z=z_*$ 时 $C=0$；
$z_*<z<1$ 时 $C>0$。膨胀经过负热容区时仍吸热，却因对外做功更多而降温。
在零点，内能减少恰好补偿膨胀功，元热量为零。$C$ 的单位始终与 $R,c_v$ 相同。
\end{xesolution}

\Needspace{10\baselineskip}
\begin{origexamplebox}{例题 20.4【恒定的热容】}
小明猜想，可能存在一些过程使过程热容保持不变，例如常见的等容、等压、等温和绝热过程。但逐一枚举
容易遗漏。能否得到“过程热容为常量”的统一过程方程？
\end{origexamplebox}

\begin{xesolution}
设一摩尔、定组成理想气体沿准静态过程的摩尔过程热容为有限常量 $C$，并取常数 $c_v,c_p$。由定义与第一定律
\begin{equation*}
  C\dd T=c_v\dd T+P\dd V,
\end{equation*}
即
\begin{equation*}
  (C-c_v)\dd T=P\dd V.
\end{equation*}
又由 $RT=PV$，
\begin{equation*}
  R\dd T=P\dd V+V\dd P.
\end{equation*}
消去 $\dd T$，得到
\begin{equation*}
  (C-c_v)V\dd P+(C-c_p)P\dd V=0.
\end{equation*}
在 $C\ne c_v$ 时两边除以 $PV(C-c_v)$ 并积分：
\begin{equation*}
  \frac{\dd P}{P}+k\frac{\dd V}{V}=0,
  \qquad
  k=\frac{C-c_p}{C-c_v},
\end{equation*}
所以统一的过程方程是
\begin{equation}
  \boxed{PV^k=\mathrm{const}},
  \qquad
  \boxed{k=\frac{C-c_p}{C-c_v}},
\end{equation}
或反解为
\begin{equation}
  \boxed{C=c_v+\frac{R}{1-k}}.
\end{equation}
这就是多方过程。它包含若干重要极限：
\begin{center}
\small
\begin{tabular}{c c c}
\toprule
过程 & 多方指数 $k$ & 过程热容 $C$\\
\midrule
等压 & $0$ & $c_p$\\
等温 & $1$ & 形式上发散（$\dd T=0$）\\
绝热 & $\gamma$ & $0$\\
等容 & $k\to\infty$ & $c_v$\\
\bottomrule
\end{tabular}
\end{center}
特别地，$C=c_v$ 时不能除以 $C-c_v$，应回到 $(C-c_v)\dd T=P\dd V$，直接得到 $\dd V=0$。
$k=1$ 的等温过程通常有 $\delta Q\ne0$、$\dd T=0$，并没有一个有限的过程热容，
只能列为 $|C|\to\infty$ 的极限。因此有限恒热容过程除等容这一单列情形外，构成
$PV^k=\mathrm{const}$ 的一整族过程；不能把“无穷大”当作与 $c_p$ 一样的有限常数。
\end{xesolution}

\Needspace{10\baselineskip}
\begin{origexamplebox}{例题 20.5【等温大气】}
小明研究一个简化的大气模型。假设大气经过长时间热接触后处处具有相同温度 $T$，气体平均摩尔质量
为 $\mu$。考虑气体密度随压强改变，求压强和分子数密度随高度 $h$ 的变化关系。
\end{origexamplebox}

\begin{xesolution}
取高度 $h$ 处厚度为 $\dd h$ 的水平气层。静力平衡要求
\begin{equation*}
  \dd P=-\rho g\,\dd h.
\end{equation*}
等温理想气体满足
\begin{equation*}
  \rho=\frac{\mu P}{RT},
\end{equation*}
故
\begin{equation*}
  \frac{\dd P}{P}=-\frac{\mu g}{RT}\dd h.
\end{equation*}
取地面 $h=0$ 处压强为 $P_0$，积分得到
\begin{equation}
  \boxed{P(h)=P_0\exp\!\left(-\frac{\mu gh}{RT}\right)}.
\end{equation}
分子数密度 $n_N=P/(k_BT)$，所以
\begin{equation}
  \boxed{n_N(h)=n_{N0}\exp\!\left(-\frac{\mu gh}{RT}\right)}.
\end{equation}

\isothermalAtmosphereDiagram
\figcap{等温大气的指数分布。标高 $H$ 决定压强和数密度衰减的快慢。}\label{fig:isothermal-atmosphere}
图\ref{fig:isothermal-atmosphere}的指数关系要求温度、成分与重力加速度近似不随高度变化。

定义等温标高
\begin{equation*}
  \boxed{H=\frac{RT}{\mu g}},
\end{equation*}
则压强、质量密度和分子数密度都按 $\e^{-h/H}$ 衰减。真实大气并不等温，重力加速度和气体成分也会随
高度改变；本题给出的是最基本的指数大气模型。
\end{xesolution}

\Needspace{10\baselineskip}
\begin{origexamplebox}{例题 20.6【毛细现象】}
小明把一根内半径为 $R$ 的玻璃管插入水面。假设水完全润湿玻璃，表面张力系数为 $\sigma$，水的密度
为 $\rho$。水面会上升多高？小明认为大树吸水、蒸腾也是同样机制，那么仅靠毛细作用给
$\SI{10}{m}$ 高处供水，需要多细的管半径？
\end{origexamplebox}

\begin{xesolution}
完全润湿时接触角 $\theta=0$。沿管内圆周，表面张力的竖直合力为
\begin{equation*}
  F_\sigma=2\pi R\sigma\cos\theta=2\pi R\sigma.
\end{equation*}
上升水柱的重力为
\begin{equation*}
  G=\rho g(\pi R^2h).
\end{equation*}
平衡时二者相等，故
\begin{equation}
  \boxed{h=\frac{2\sigma}{\rho gR}}.
\end{equation}

\capillaryRevisionDiagram
\figcap{完全润湿的毛细管。弯月面中心低、壁边高，壁边表面张力沿切线向上。}\label{fig:capillary-wetting}
图\ref{fig:capillary-wetting}的 $h$ 从外部水平水面量起。这里忽略弯月面附近体积相对长水柱的微小修正；
接触角为零使周界张力全用于托住水柱，量纲也给出 $[\sigma/(\rho gR)]=\mathrm m$。

若要求 $h=\SI{10}{m}$，则
\begin{equation*}
  \boxed{R=\frac{2\sigma}{\rho gh}}.
\end{equation*}
取室温水 $\sigma\approx\SI{0.073}{N.m^{-1}}$、$\rho\approx\SI{1.0e3}{kg.m^{-3}}$，得到
\begin{equation*}
  R\approx\frac{2\times0.073}{10^3\times9.8\times10}
  \approx\SI{1.5e-6}{m}.
\end{equation*}
这说明微米级管道的毛细压强足以对应十米水柱。但真实树木的远距离输水并非单靠静态毛细平衡：叶片
蒸腾在木质部水柱中产生负压，水分子的内聚力和导管结构共同维持连续水柱；毛细作用主要帮助润湿和
稳定细小导管。
\end{xesolution}

\Needspace{10\baselineskip}
\begin{origexamplebox}{例题 20.7【吹泡泡】}
小明拿出泡泡水吹泡泡。设泡泡水的表面张力系数为 $\sigma$，泡泡为半径 $R$ 的正球。

（1）若不计液膜厚度，泡泡内部压强为多大？

（2）若平衡时液膜的质量面密度为 $\sigma_0$，忽略空气质量，把空气看成双原子理想气体。在小的径向
扰动下，设气泡仍能维持球形，判断平衡是否稳定。
\end{origexamplebox}

\begin{xesolution}
（1）肥皂泡有内、外两个气液界面，每个球面都产生 $2\sigma/R$ 的 Laplace 压强差，因此总压强差为
\begin{equation*}
  p_{\rm in}-p_0=\frac{4\sigma}{R}.
\end{equation*}
所以
\begin{equation}
  \boxed{p_{\rm in}=p_0+\frac{4\sigma}{R}}.
\end{equation}
注意：只有一个界面的液滴压强差是 $2\sigma/R$；肥皂泡有两个界面，不能漏掉因子2。

\bubbleRevisionDiagram
\figcap{肥皂泡的两界面与径向扰动。$p_R$ 是平衡时内部压强；虚线表示均匀向外位移。}\label{fig:bubble-radius}
也可直接对图\ref{fig:bubble-radius}作虚功计算：$\dd V=4\pi r^2\dd r$，
两界面的面积增量为 $\dd A_s=16\pi r\dd r$，平衡要求
$(p_R-p_0)\dd V=\sigma\dd A_s$，由此再次得到因子4。

（2）设半径受扰动变为 $r=R+x$，$|x|\ll R$。小振幅并不自动意味着变化很快：
若传热足够快，气体接近等温；若振动周期远短于漏热时间、又足以维持内部压强近似均匀，则接近绝热。
先统一用多方指数 $\kappa$ 表示这两种极限：等温为1，常温双原子气体的可逆绝热近似为 $7/5$。
取表面张力 $\sigma$ 在扰动中不变，忽略气体惯性、膜的黏性耗散和外部空气附加质量。记 $p_R=p_{\rm in}$，于是
\begin{equation*}
  p(r)r^{3\kappa}=p_R R^{3\kappa},
  \qquad
  p'(R)=-\frac{3\kappa p_R}{R}.
\end{equation*}
球面对称时所有真实力的矢量合力为零，不能把它作为整个泡的平移合力。这里计算的是与半径 $r$
共轭的广义力，即虚功中 $\dd r$ 的系数：压强贡献为 $4\pi r^2[p(r)-p_0]$。
两个界面的表面张力功对应有效势 $E_s=8\pi\sigma r^2$，向内的广义力为
\begin{equation*}
  -\frac{\dd E_s}{\dd r}=-16\pi\sigma r.
\end{equation*}
故总径向广义力为
\begin{equation*}
  F(r)=4\pi r^2[p(r)-p_0]-16\pi\sigma r.
\end{equation*}
在 $r=R$ 处利用平衡条件 $p_{\rm in}-p_0=4\sigma/R$，得到
\begin{equation*}
  F'(R)=4\pi\bigl(4\sigma-3\kappa Rp_R\bigr).
\end{equation*}
液膜质量为 $m_f=4\pi R^2\sigma_0$。线性化运动方程写成
\begin{equation*}
  m_f\ddot x=F'(R)x=-m_f\omega^2x,
\end{equation*}
其中
\begin{equation}
  \boxed{\omega^2=
  \frac{3\kappa Rp_R-4\sigma}{R^2\sigma_0}}.
\end{equation}
代入 $p_R=p_0+4\sigma/R$，分子为
\begin{equation*}
  3\kappa Rp_0+(12\kappa-4)\sigma>0\qquad(\kappa=1\ \text{或}\ 7/5).
\end{equation*}
所以 $\omega^2>0$，小的球对称径向扰动会产生回复力，
\begin{equation*}
  \boxed{\text{该平衡对题设的小径向扰动稳定。}}
\end{equation*}
这里判断的是理想化的球对称模式；真实泡膜还可能因排液、蒸发、膜厚不均或非球形扰动而破裂。
由 $[\sigma]=\mathrm{N/m}$、$[\sigma_0]=\mathrm{kg/m^2}$ 可验得 $[\omega^2]=\mathrm{s^{-2}}$。
一般传热速率介于两极限之间时，还可能出现相位滞后和阻尼，题设不足以指定唯一振动频率；但上述两极限都稳定。
\end{xesolution}

\Needspace{8\baselineskip}
\begin{origexamplebox}{例题 20.8【华氏温标】}
美国常用的华氏温标把标准大气压下冰的熔点定为 $32\,{}^\circ\mathrm F$，水的沸点定为
$212\,{}^\circ\mathrm F$。

（1）导出华氏温度与摄氏温度之间的换算关系；（2）给出绝对零度的华氏温度。
\end{origexamplebox}

\begin{xesolution}
摄氏温标中冰点与沸点相差 $100$ 度，华氏温标中相差 $212-32=180$ 度，因此同一温差满足
\begin{equation*}
  \Delta t_F=\frac{180}{100}\Delta t_C=\frac95\Delta t_C.
\end{equation*}
再用 $t_C=0$ 对应 $t_F=32$，得到
\begin{equation}
  \boxed{t_F=\frac95t_C+32},
  \qquad
  \boxed{t_C=\frac59(t_F-32)}.
\end{equation}
绝对零度 $t_C=-273.15\,{}^\circ\mathrm C$，所以
\begin{equation*}
  t_F=\frac95(-273.15)+32=-459.67\,{}^\circ\mathrm F.
\end{equation*}
即
\begin{equation}
  \boxed{0\,\mathrm K=-459.67\,{}^\circ\mathrm F}.
\end{equation}
温差换算只乘 $9/5$，绝对温度读数换算还必须处理零点偏移；这是两类公式最容易混淆之处。
\end{xesolution}

\Needspace{11\baselineskip}
\begin{origexamplebox}{例题 20.9【循环过程】}
如图所示，一定量单原子理想气体从初态 $a(P_1,V_1)$ 出发，先经等容过程到状态 $b$，再经等压过程
到状态 $c$，最后经等温过程回到 $a$，完成一个循环。图中 $P_b=P_c=P_1/4$。已知一摩尔气体温度
改变 $\SI{1}{K}$ 时内能改变 $c_v=3R/2$。

（1）$c\to a$ 吸热还是放热，热量多少？（2）$a\to b$ 吸热还是放热，热量多少？

（3）$b\to c$ 吸热还是放热，热量多少？（4）整个循环的净功和净吸放热量各为多少？
\cycleQuestionRedraw
\diagramcaption{图 20.9：$a\to b\to c\to a$ 循环。箭头给出过程方向，$c\to a$ 为等温压缩。}
图20.9中 $a,c$ 在同一条双曲线上，而 $b,c$ 具有相同压强。
\end{origexamplebox}

\begin{xesolution}
过程 $c\to a$ 等温，所以 $P_cV_c=P_aV_a$。由 $P_c=P_1/4$ 得
\begin{equation*}
  \frac{P_1}{4}V_c=P_1V_1,
  \qquad
  \boxed{V_c=4V_1}.
\end{equation*}
设气体物质的量为 $n$，并统一取 $Q>0$ 为吸热、$W^{\rm out}>0$ 为气体对外做功。

\cycleExampleDiagram
\figcap{循环计算所用状态。$a,b$ 等体积，$b,c$ 等压，$a,c$ 等温。}\label{fig:cycle-solution}
以下逐段热量与图\ref{fig:cycle-solution}的三条有向路径对应。

（1）$c\to a$ 为温度 $T_a=P_1V_1/(nR)$ 的等温压缩，$\Delta U_{ca}=0$，
\begin{equation*}
  Q_{ca}=W^{\rm out}_{ca}
  =nRT_a\ln\frac{V_a}{V_c}
  =P_1V_1\ln\frac14.
\end{equation*}
因此
\begin{equation}
  \boxed{Q_{ca}=-P_1V_1\ln4},
\end{equation}
即气体放热 $P_1V_1\ln4$。

（2）$a\to b$ 等容，且 $T_b=T_a/4$。所以
\begin{equation*}
  W^{\rm out}_{ab}=0,
  \qquad
  Q_{ab}=\Delta U_{ab}
  =\frac32nR\left(\frac{T_a}{4}-T_a\right)
  =-\frac98P_1V_1.
\end{equation*}
故
\begin{equation}
  \boxed{Q_{ab}=-\frac98P_1V_1},
\end{equation}
即放热 $9P_1V_1/8$。

（3）$b\to c$ 等压，温度由 $T_a/4$ 升到 $T_a$。单原子气体 $c_p=c_v+R=5R/2$，所以
\begin{equation*}
  Q_{bc}=nc_p(T_c-T_b)
  =\frac52nR\cdot\frac34T_a
  =\frac{15}{8}P_1V_1.
\end{equation*}
即
\begin{equation}
  \boxed{Q_{bc}=\frac{15}{8}P_1V_1},
\end{equation}
气体吸热。

（4）净功可由各段相加。等容段功为零，等压段功为
\begin{equation*}
  W^{\rm out}_{bc}=\frac{P_1}{4}(4V_1-V_1)=\frac34P_1V_1,
\end{equation*}
等温压缩段功为 $-P_1V_1\ln4$。因此
\begin{equation}
  \boxed{W^{\rm out}_{\rm cyc}=P_1V_1\left(\frac34-\ln4\right)<0}.
\end{equation}
循环内能改变量为零，故净热量与净功相等：
\begin{equation}
  \boxed{Q_{\rm cyc}=P_1V_1\left(\frac34-\ln4\right)<0}.
\end{equation}
所以该循环是逆时针循环，外界对气体做净功
$P_1V_1(\ln4-3/4)$，气体向外净放出同样多的热量。
\end{xesolution}

\Needspace{11\baselineskip}
\begin{origexamplebox}{例题 20.10【直线过程】}
如图所示，一摩尔理想气体从 $P$--$V$ 图中的状态 $A$ 出发，沿缓慢直线过程到达状态 $B$。已知
$P_B/P_A=1/2$。若要使全过程的最终结果是气体从外界吸收热量，状态 $B$ 与状态 $A$ 的体积比应满足
什么条件？已知该理想气体每摩尔内能为 $U=3RT/2$。
\straightQuestionRedraw
\diagramcaption{图 20.10：状态 $A$ 到状态 $B$ 的准静态直线过程；解答中的详图进一步补出端点压强、体积及面积关系。}
图20.10中 $B$ 的纵坐标为 $A$ 的一半，这个比例将同时用于内能变化和梯形面积。
\end{origexamplebox}

\begin{xesolution}
令
\begin{equation*}
  x=\frac{V_B}{V_A},
  \qquad P_B=\frac12P_A.
\end{equation*}
直线过程的对外功等于图下梯形面积：
\begin{equation*}
  W^{\rm out}
  =\frac{P_A+P_B}{2}(V_B-V_A)
  =\frac34P_AV_A(x-1).
\end{equation*}

\straightProcessDiagram
\figcap{直线过程的两个端点及轴上投影。积分区间是 $V_A$ 到 $V_B$。}\label{fig:straight-solution}
图\ref{fig:straight-solution}中直线下方的梯形面积给功，端点的 $PV$ 差给内能变化，二者要分别计算。

由 $U=3RT/2$ 以及一摩尔理想气体 $RT=PV$，
\begin{align*}
  \Delta U
  &=\frac32(P_BV_B-P_AV_A)\\
  &=\frac34P_AV_A(x-2).
\end{align*}
因此吸热量
\begin{align*}
  Q&=\Delta U+W^{\rm out}\\
   &=\frac34P_AV_A[(x-2)+(x-1)]\\
   &=\frac34P_AV_A(2x-3).
\end{align*}
要求气体净吸热即 $Q>0$，得到
\begin{equation}
  \boxed{\frac{V_B}{V_A}>\frac32}.
\end{equation}
若 $1<x<3/2$，气体虽然膨胀并对外做功，但温度、内能下降得更多，仍会净放热；因此“膨胀”不等于
“吸热”。
\end{xesolution}

\Needspace{11\baselineskip}
\begin{origexamplebox}{例题 20.11【提升活塞】}
如图所示，在质量为 $M$、内部横截面积为 $A$ 的竖直绝热气缸中，活塞封闭了一定量初温为 $T_0$ 的
理想气体。活塞也绝热，活塞质量及活塞与气缸之间的摩擦均可忽略。大气压强为 $p_0$，重力加速度为
$g$。气缸由不密封的小支脚支承，大气能作用于缸底外侧；忽略气体质量，取恒定绝热指数
$\gamma=c_p/c_v$，并满足 $Mg<p_0A$。现把活塞缓慢上提；当活塞到达气缸开口处、气体尚未泄出时，
气缸恰好离开支承面。求此时气体温度。
\liftingPistonQuestionRedraw
\diagramcaption{图 20.11：初态与刚离地终态。终态取活塞仍密封气体的极限位置；缸底外侧受到大气压。}
图20.11右侧取开口前最后密封的状态；活塞再向上、气体泄出后就不能继续使用同一封闭系统的绝热方程。
\end{origexamplebox}

\begin{xesolution}
初始时活塞没有质量，又未受额外拉力，故气体压强为 $p_i=p_0$。上提过程缓慢，使每一时刻都可视为
平衡态；气缸、活塞又绝热，所以封闭气体经历可逆绝热膨胀。

设某时刻气体压强为 $p$，外界对活塞的上拉力为 $F$。质量可忽略的活塞受力平衡给出
\begin{equation*}
  F+pA=p_0A,
  \qquad
  F=(p_0-p)A.
\end{equation*}
还必须单独检查缸体，不能说低压气体把缸底“向上拉”。缸底内侧气体给缸体向下的力 $pA$，
外侧大气给向上的力 $p_0A$，支承力为 $N$。竖直平衡给出
\[
 N+p_0A=pA+Mg,\qquad N=Mg-(p_0-p)A.
\]
随着气体膨胀，$p$ 下降，支承力逐渐减小。刚离地时 $N=0$，于是
\begin{equation*}
  \boxed{p_f=p_0-\frac{Mg}{A}}.
\end{equation*}

\liftingPistonDiagram
\figcap{活塞与缸体分别受力。力箭头的作用对象不同，不能直接在一张受力图中混加。}\label{fig:lift-forces}
图\ref{fig:lift-forces}左图给 $F=(p_0-p)A$，右图在 $N=0$ 时给 $(p_0-p_f)A=Mg$，因此才有 $F_f=Mg$。

由绝热关系 $T^\gamma/P^{\gamma-1}=\mathrm{const}$，
\begin{equation*}
  \frac{T_f^\gamma}{p_f^{\gamma-1}}
  =\frac{T_0^\gamma}{p_0^{\gamma-1}}.
\end{equation*}
因此
\begin{equation}
  \boxed{T_f=T_0\left(\frac{p_f}{p_0}\right)^{(\gamma-1)/\gamma}
  =T_0\left(1-\frac{Mg}{p_0A}\right)^{(\gamma-1)/\gamma}}.
\end{equation}
要使末态压强为正并能准静态达到，题设参数必须满足 $Mg<p_0A$。上提过程中气体膨胀并对外做功，却
没有从环境吸热，因此内能和温度都降低，结果中的括号确实小于1。
\end{xesolution}

